\documentclass[12pt,a4paper]{article}

% Packages
\usepackage[utf8]{inputenc}
\usepackage[margin=1in]{geometry}
\usepackage{graphicx}
\usepackage{float}
\usepackage{hyperref}
\usepackage{listings}
\usepackage{xcolor}
\usepackage{titlesec}
\usepackage{fancyhdr}
\usepackage{tabularx}
\usepackage{longtable}
\usepackage{array}
\usepackage{amsmath}
\usepackage{amssymb}
\usepackage{booktabs}
\usepackage{multirow}
\usepackage{enumitem}

% Page setup
\pagestyle{fancy}
\fancyhf{}
\rhead{SDS - Expert-Guided Multimodal AI Ecosystem}
\lhead{SEECS, NUST}
\rfoot{Page \thepage}

% Title formatting
\titleformat{\section}
  {\normalfont\Large\bfseries}{\thesection}{1em}{}
\titleformat{\subsection}
  {\normalfont\large\bfseries}{\thesubsection}{1em}{}

% Code listing settings
\lstset{
    basicstyle=\ttfamily\small,
    commentstyle=\color{gray},
    keywordstyle=\color{blue},
    stringstyle=\color{red},
    showstringspaces=false,
    breaklines=true,
    frame=single,
    numbers=left,
    numberstyle=\tiny\color{gray}
}

\begin{document}

% Title Page
\begin{titlepage}
    \centering
    \vspace*{2cm}
    
    {\huge\bfseries Software Design Specifications\par}
    \vspace{1cm}
    
    {\Large\bfseries An Expert-Guided Multimodal AI Ecosystem for\par}
    {\Large\bfseries Brain MRI Analysis and Clinical Report Generation\par}
    \vspace{2cm}
    
    {\Large Final Year Design Project\par}
    \vspace{1cm}
    
    {\large School of Electrical Engineering and Computer Science (SEECS)\par}
    {\large National University of Sciences and Technology (NUST)\par}
    \vspace{2cm}
    
    {\large\bfseries Submitted By:\par}
    {\large Farhan Ahmed\par}
    \vspace{1cm}
    
    {\large\bfseries Submission Date:\par}
    {\large \today\par}
    
    \vfill
\end{titlepage}

% Table of Contents
\tableofcontents
\newpage

% Document History
\section*{Document History}
\addcontentsline{toc}{section}{Document History}

\begin{table}[H]
\centering
\begin{tabularx}{\textwidth}{|c|c|X|c|}
\hline
\textbf{Version} & \textbf{Date} & \textbf{Description} & \textbf{Author} \\
\hline
1.0 & \today & Initial SDS Document & Farhan Ahmed \\
\hline
\end{tabularx}
\end{table}

\newpage

% 1. Introduction
\section{Introduction}

\subsection{Purpose}
This Software Design Specifications (SDS) document provides a comprehensive architectural and detailed design for the Expert-Guided Multimodal AI Ecosystem for Brain MRI Analysis and Clinical Report Generation. The document serves as a blueprint for the implementation, testing, and deployment of the system.

\subsection{Scope}
The system is designed to provide an end-to-end solution for medical image analysis with the following key capabilities:

\begin{itemize}
    \item Multi-modal brain MRI segmentation using Mixture of Modality Experts (MoME+) architecture
    \item Continual learning capabilities with knowledge retention across multiple datasets
    \item Brain atlas mapping for anatomical localization
    \item Automated clinical report generation using fine-tuned Large Language Models
    \item Web-based dashboard for visualization and interaction
    \item RESTful API for integration with clinical workflows
\end{itemize}

\subsection{Definitions, Acronyms, and Abbreviations}

\begin{table}[H]
\centering
\begin{tabularx}{\textwidth}{|l|X|}
\hline
\textbf{Term} & \textbf{Definition} \\
\hline
MoME+ & Mixture of Modality Experts Plus - Enhanced architecture for multi-modal medical image segmentation \\
\hline
BraTS & Brain Tumor Segmentation Challenge Dataset \\
\hline
FLAIR & Fluid-Attenuated Inversion Recovery MRI sequence \\
\hline
WT & Whole Tumor region \\
\hline
TC & Tumor Core region \\
\hline
ET & Enhancing Tumor region \\
\hline
EWC & Elastic Weight Consolidation - Continual learning technique \\
\hline
LoRA & Low-Rank Adaptation for efficient LLM fine-tuning \\
\hline
PEFT & Parameter-Efficient Fine-Tuning \\
\hline
CBAM & Convolutional Block Attention Module \\
\hline
SE & Squeeze-and-Excitation attention mechanism \\
\hline
MNI152 & Montreal Neurological Institute standard brain atlas \\
\hline
API & Application Programming Interface \\
\hline
REST & Representational State Transfer \\
\hline
LLM & Large Language Model \\
\hline
\end{tabularx}
\end{table}

\subsection{References}
\begin{enumerate}
    \item Software Requirements Specification (SRS) Document - Project SRS
    \item IEEE Std 1016-2009 - IEEE Standard for Information Technology Systems Design
    \item BraTS Challenge Documentation
    \item MedGemma-4B Model Documentation
    \item Django REST Framework Documentation
\end{enumerate}

\subsection{Overview}
This document is organized into the following sections:

\begin{itemize}
    \item \textbf{Section 2:} System Architecture - High-level architectural design
    \item \textbf{Section 3:} Data Design - Database schema and data structures
    \item \textbf{Section 4:} Component Design - Detailed module specifications
    \item \textbf{Section 5:} Interface Design - API and UI specifications
    \item \textbf{Section 6:} Security Design - Authentication and authorization
    \item \textbf{Section 7:} Deployment Design - Infrastructure and deployment
\end{itemize}

\newpage

% 2. System Architecture
\section{System Architecture}

\subsection{Architectural Overview}
The system follows a modular, microservices-inspired architecture with clear separation of concerns. The architecture is designed to support scalability, maintainability, and extensibility.

\begin{figure}[H]
\centering
\includegraphics[width=\textwidth]{diagrams/system_architecture.png}
\caption{High-Level System Architecture}
\label{fig:system_arch}
\end{figure}

\textit{Note: Generate diagram using PlantUML code provided in diagrams/system\_architecture.puml}

\subsection{Architectural Patterns}

\subsubsection{Layered Architecture}
The system employs a 4-tier layered architecture:

\begin{enumerate}
    \item \textbf{Presentation Layer:} Web-based dashboard (React/TypeScript)
    \item \textbf{Application Layer:} Business logic and orchestration (Django REST Framework)
    \item \textbf{Domain Layer:} Core ML models and algorithms (PyTorch)
    \item \textbf{Data Layer:} Data persistence and management (PostgreSQL, File Storage)
\end{enumerate}

\subsubsection{Model-View-Controller (MVC)}
The backend follows the MVC pattern:
\begin{itemize}
    \item \textbf{Model:} Django ORM models for data representation
    \item \textbf{View:} DRF ViewSets for request handling
    \item \textbf{Controller:} Business logic in service classes
\end{itemize}

\subsubsection{Pipeline Pattern}
The inference workflow follows a pipeline pattern:
\begin{equation}
\text{Input} \xrightarrow{\text{Preprocess}} \text{Segmentation} \xrightarrow{\text{Atlas Mapping}} \text{JSON} \xrightarrow{\text{LLM}} \text{Report}
\end{equation}

\subsection{Component Diagram}

\begin{figure}[H]
\centering
\includegraphics[width=0.9\textwidth]{diagrams/component_diagram.png}
\caption{Component Diagram}
\label{fig:component}
\end{figure}

\textit{Note: Generate diagram using PlantUML code provided in diagrams/component\_diagram.puml}

\subsection{Technology Stack}

\begin{table}[H]
\centering
\begin{tabularx}{\textwidth}{|l|l|X|}
\hline
\textbf{Layer} & \textbf{Technology} & \textbf{Purpose} \\
\hline
Frontend & React 18+ & User interface framework \\
 & TypeScript & Type-safe development \\
 & TailwindCSS & Styling framework \\
\hline
Backend & Django 4.2+ & Web framework \\
 & Django REST Framework & API development \\
 & Celery & Asynchronous task processing \\
 & Redis & Message broker and caching \\
\hline
ML Framework & PyTorch 2.0+ & Deep learning framework \\
 & MONAI & Medical imaging library \\
 & Transformers & LLM integration \\
 & nibabel & NIfTI file handling \\
\hline
Database & PostgreSQL 14+ & Relational database \\
 & MinIO / S3 & Object storage for MRI files \\
\hline
Deployment & Docker & Containerization \\
 & Docker Compose & Multi-container orchestration \\
 & Nginx & Reverse proxy \\
 & Gunicorn & WSGI server \\
\hline
\end{tabularx}
\end{table}

\newpage

% 3. Data Design
\section{Data Design}

\subsection{Entity-Relationship Diagram}

\begin{figure}[H]
\centering
\includegraphics[width=\textwidth]{diagrams/erd_diagram.png}
\caption{Entity-Relationship Diagram}
\label{fig:erd}
\end{figure}

\textit{Note: Generate diagram using PlantUML code provided in diagrams/erd\_diagram.puml}

\subsection{Database Schema}

\subsubsection{Patient Table}
\begin{table}[H]
\centering
\begin{tabularx}{\textwidth}{|l|l|l|X|}
\hline
\textbf{Field} & \textbf{Type} & \textbf{Constraints} & \textbf{Description} \\
\hline
patient\_id & UUID & PRIMARY KEY & Unique patient identifier \\
patient\_code & VARCHAR(50) & UNIQUE, NOT NULL & Anonymized patient code \\
age & INTEGER & CHECK (age > 0) & Patient age \\
sex & CHAR(1) & CHECK (sex IN ('M','F','O')) & Patient sex \\
created\_at & TIMESTAMP & DEFAULT NOW() & Record creation time \\
updated\_at & TIMESTAMP & ON UPDATE NOW() & Last update time \\
\hline
\end{tabularx}
\end{table}

\subsubsection{MRI Scan Table}
\begin{table}[H]
\centering
\begin{tabularx}{\textwidth}{|l|l|l|X|}
\hline
\textbf{Field} & \textbf{Type} & \textbf{Constraints} & \textbf{Description} \\
\hline
scan\_id & UUID & PRIMARY KEY & Unique scan identifier \\
patient\_id & UUID & FOREIGN KEY & Reference to patient \\
scan\_date & DATE & NOT NULL & Date of MRI acquisition \\
t1\_path & VARCHAR(500) & NOT NULL & Path to T1 modality \\
t1ce\_path & VARCHAR(500) & NOT NULL & Path to T1ce modality \\
t2\_path & VARCHAR(500) & NOT NULL & Path to T2 modality \\
flair\_path & VARCHAR(500) & NOT NULL & Path to FLAIR modality \\
status & VARCHAR(20) & DEFAULT 'uploaded' & Processing status \\
created\_at & TIMESTAMP & DEFAULT NOW() & Upload timestamp \\
\hline
\end{tabularx}
\end{table}

\subsubsection{Segmentation Result Table}
\begin{table}[H]
\centering
\begin{tabularx}{\textwidth}{|l|l|l|X|}
\hline
\textbf{Field} & \textbf{Type} & \textbf{Constraints} & \textbf{Description} \\
\hline
seg\_id & UUID & PRIMARY KEY & Unique segmentation ID \\
scan\_id & UUID & FOREIGN KEY & Reference to MRI scan \\
model\_version & VARCHAR(50) & NOT NULL & Model version used \\
mask\_path & VARCHAR(500) & NOT NULL & Path to segmentation mask \\
wt\_volume & FLOAT & CHECK (>= 0) & Whole tumor volume (cm³) \\
tc\_volume & FLOAT & CHECK (>= 0) & Tumor core volume (cm³) \\
et\_volume & FLOAT & CHECK (>= 0) & Enhancing tumor volume (cm³) \\
dice\_wt & FLOAT & CHECK (0 <= dice\_wt <= 1) & Dice score for WT \\
dice\_tc & FLOAT & CHECK (0 <= dice\_tc <= 1) & Dice score for TC \\
dice\_et & FLOAT & CHECK (0 <= dice\_et <= 1) & Dice score for ET \\
inference\_time & FLOAT & & Time taken (seconds) \\
created\_at & TIMESTAMP & DEFAULT NOW() & Inference timestamp \\
\hline
\end{tabularx}
\end{table}

\subsubsection{Atlas Mapping Table}
\begin{table}[H]
\centering
\begin{tabularx}{\textwidth}{|l|l|l|X|}
\hline
\textbf{Field} & \textbf{Type} & \textbf{Constraints} & \textbf{Description} \\
\hline
mapping\_id & UUID & PRIMARY KEY & Unique mapping ID \\
seg\_id & UUID & FOREIGN KEY & Reference to segmentation \\
atlas\_type & VARCHAR(50) & NOT NULL & Atlas used (e.g., Harvard-Oxford) \\
affected\_regions & JSONB & NOT NULL & Array of affected brain regions \\
alignment\_quality & FLOAT & CHECK (0 <= value <= 1) & Registration quality score \\
created\_at & TIMESTAMP & DEFAULT NOW() & Mapping timestamp \\
\hline
\end{tabularx}
\end{table}

\subsubsection{Clinical Report Table}
\begin{table}[H]
\centering
\begin{tabularx}{\textwidth}{|l|l|l|X|}
\hline
\textbf{Field} & \textbf{Type} & \textbf{Constraints} & \textbf{Description} \\
\hline
report\_id & UUID & PRIMARY KEY & Unique report ID \\
seg\_id & UUID & FOREIGN KEY & Reference to segmentation \\
json\_descriptor & JSONB & NOT NULL & Structured JSON descriptor \\
report\_text & TEXT & NOT NULL & Generated clinical report \\
llm\_model\_version & VARCHAR(50) & NOT NULL & LLM version used \\
confidence\_score & FLOAT & CHECK (0 <= value <= 1) & Report confidence \\
reviewed & BOOLEAN & DEFAULT FALSE & Clinician review status \\
reviewer\_feedback & TEXT & & Clinician comments \\
created\_at & TIMESTAMP & DEFAULT NOW() & Generation timestamp \\
reviewed\_at & TIMESTAMP & & Review timestamp \\
\hline
\end{tabularx}
\end{table}

\subsubsection{User Table}
\begin{table}[H]
\centering
\begin{tabularx}{\textwidth}{|l|l|l|X|}
\hline
\textbf{Field} & \textbf{Type} & \textbf{Constraints} & \textbf{Description} \\
\hline
user\_id & UUID & PRIMARY KEY & Unique user ID \\
username & VARCHAR(100) & UNIQUE, NOT NULL & Login username \\
email & VARCHAR(255) & UNIQUE, NOT NULL & User email \\
password\_hash & VARCHAR(255) & NOT NULL & Hashed password \\
role & VARCHAR(20) & NOT NULL & User role (admin/radiologist/researcher) \\
is\_active & BOOLEAN & DEFAULT TRUE & Account status \\
last\_login & TIMESTAMP & & Last login time \\
created\_at & TIMESTAMP & DEFAULT NOW() & Account creation \\
\hline
\end{tabularx}
\end{table}

\subsection{Data Flow Diagram}

\begin{figure}[H]
\centering
\includegraphics[width=\textwidth]{diagrams/data_flow_diagram.png}
\caption{Level 0 Data Flow Diagram}
\label{fig:dfd0}
\end{figure}

\textit{Note: Generate diagram using PlantUML code provided in diagrams/data\_flow\_diagram.puml}

\subsection{JSON Schema Design}

The system uses a structured JSON schema for anatomical descriptors:

\begin{lstlisting}[language=json]
{
  "case_id": "string",
  "patient_metadata": {
    "age": "integer",
    "sex": "string ('M'/'F'/'O')"
  },
  "segmentation_analysis": {
    "whole_tumor": {
      "volume_cm3": "float",
      "voxel_count": "integer"
    },
    "tumor_core": { ... },
    "enhancing_tumor": { ... }
  },
  "anatomical_localization": {
    "affected_regions": [
      {
        "region_name": "string",
        "hemisphere": "string",
        "involvement_percentage": "float",
        "volume_cm3": "float"
      }
    ]
  },
  "tumor_characteristics": {
    "enhancing_component": "boolean",
    "necrotic_component": "boolean",
    "edema_present": "boolean"
  },
  "model_metadata": {
    "model_version": "string",
    "dice_scores": { ... },
    "inference_timestamp": "ISO8601"
  }
}
\end{lstlisting}

\newpage

% 4. Component Design
\section{Component Design}

\subsection{Preprocessing Module}

\subsubsection{Purpose}
Standardize and prepare multi-modal MRI data for model inference.

\subsubsection{Class Diagram}

\begin{figure}[H]
\centering
\includegraphics[width=0.9\textwidth]{diagrams/preprocessing_class.png}
\caption{Preprocessing Module Class Diagram}
\label{fig:preprocess_class}
\end{figure}

\textit{Note: Generate diagram using PlantUML code provided in diagrams/preprocessing\_class.puml}

\subsubsection{Key Algorithms}

\paragraph{Intensity Normalization}
\begin{equation}
I_{norm} = \frac{I - \mu}{\sigma}
\end{equation}
where $I$ is the input intensity, $\mu$ is the mean, and $\sigma$ is the standard deviation.

\paragraph{Spatial Resampling}
Data is resampled to isotropic resolution (1mm³) using trilinear interpolation:
\begin{equation}
V_{resampled}(x,y,z) = \sum_{i,j,k} V_{original}(i,j,k) \cdot w(x-i) \cdot w(y-j) \cdot w(z-k)
\end{equation}

\subsection{MoME+ Segmentation Module}

\subsubsection{Architecture}

The MoME+ model consists of:
\begin{itemize}
    \item 4 modality-specific expert networks (3D U-Nets)
    \item Hierarchical gating network with attention
    \item Multi-scale feature fusion
\end{itemize}

\begin{figure}[H]
\centering
\includegraphics[width=\textwidth]{diagrams/mome_architecture.png}
\caption{MoME+ Architecture Diagram}
\label{fig:mome_arch}
\end{figure}

\textit{Note: Generate diagram using PlantUML code provided in diagrams/mome\_architecture.puml}

\subsubsection{Mathematical Formulation}

\paragraph{Expert Network Output}
For each modality $m \in \{T1, T1ce, T2, FLAIR\}$:
\begin{equation}
F_m = Expert_m(X_m, \theta_m)
\end{equation}
where $X_m$ is the input modality and $\theta_m$ are expert-specific parameters.

\paragraph{Gating Mechanism}
\begin{equation}
G = \sigma(Conv_{gate}([F_{T1}, F_{T1ce}, F_{T2}, F_{FLAIR}]))
\end{equation}

\paragraph{Feature Fusion}
\begin{equation}
F_{fused} = \sum_{m=1}^{4} G_m \odot F_m
\end{equation}

\paragraph{Final Prediction}
\begin{equation}
P = Decoder(F_{fused})
\end{equation}
where $P \in \mathbb{R}^{H \times W \times D \times 4}$ for {background, WT, TC, ET}.

\subsection{Continual Learning Module}

\subsubsection{Elastic Weight Consolidation (EWC)}

The EWC loss function:
\begin{equation}
\mathcal{L}_{EWC} = \mathcal{L}_{task}(\theta) + \frac{\lambda}{2} \sum_i F_i(\theta_i - \theta^*_i)^2
\end{equation}
where:
\begin{itemize}
    \item $\mathcal{L}_{task}$ is the task-specific loss
    \item $F_i$ is the Fisher information matrix
    \item $\theta^*_i$ are the optimal parameters from the previous task
    \item $\lambda$ is the importance weight
\end{itemize}

\subsubsection{Fisher Information Computation}
\begin{equation}
F_i = \mathbb{E}_{x \sim D}\left[\left(\frac{\partial \log p(y|x,\theta)}{\partial \theta_i}\right)^2\right]
\end{equation}

\subsection{Atlas Mapping Module}

\subsubsection{Sequence Diagram}

\begin{figure}[H]
\centering
\includegraphics[width=\textwidth]{diagrams/atlas_mapping_sequence.png}
\caption{Atlas Mapping Process Sequence Diagram}
\label{fig:atlas_seq}
\end{figure}

\textit{Note: Generate diagram using PlantUML code provided in diagrams/atlas\_mapping\_sequence.puml}

\subsubsection{Registration Algorithm}

\paragraph{Affine Registration}
Optimize the transformation matrix $T$:
\begin{equation}
T^* = \arg\min_T \sum_{i} || I_{fixed}(T(x_i)) - I_{moving}(x_i) ||^2
\end{equation}

\paragraph{Overlap Calculation}
For each brain region $R$:
\begin{equation}
Overlap(R) = \frac{|Mask \cap R|}{|R|} \times 100\%
\end{equation}

\subsection{LLM Report Generation Module}

\subsubsection{Activity Diagram}

\begin{figure}[H]
\centering
\includegraphics[width=0.8\textwidth]{diagrams/report_generation_activity.png}
\caption{Report Generation Activity Diagram}
\label{fig:report_activity}
\end{figure}

\textit{Note: Generate diagram using PlantUML code provided in diagrams/report\_generation\_activity.puml}

\subsubsection{LoRA Fine-tuning}

The LoRA approach modifies the weight update:
\begin{equation}
W' = W_0 + \Delta W = W_0 + BA
\end{equation}
where:
\begin{itemize}
    \item $W_0 \in \mathbb{R}^{d \times k}$ is the frozen pre-trained weight
    \item $B \in \mathbb{R}^{d \times r}$ and $A \in \mathbb{R}^{r \times k}$ are trainable low-rank matrices
    \item $r \ll \min(d,k)$ is the rank
\end{itemize}

\paragraph{Parameter Efficiency}
Number of trainable parameters:
\begin{equation}
N_{trainable} = r \cdot (d + k) \ll d \cdot k = N_{original}
\end{equation}

\subsection{API Module}

\subsubsection{State Diagram}

\begin{figure}[H]
\centering
\includegraphics[width=0.9\textwidth]{diagrams/scan_state_diagram.png}
\caption{MRI Scan Processing State Diagram}
\label{fig:state}
\end{figure}

\textit{Note: Generate diagram using PlantUML code provided in diagrams/scan\_state\_diagram.puml}

\newpage

% 5. Interface Design
\section{Interface Design}

\subsection{API Interface Specifications}

\subsubsection{RESTful API Endpoints}

\begin{longtable}{|p{3cm}|p{2cm}|p{4cm}|p{4cm}|}
\hline
\textbf{Endpoint} & \textbf{Method} & \textbf{Request} & \textbf{Response} \\
\hline
\endfirsthead
\hline
\textbf{Endpoint} & \textbf{Method} & \textbf{Request} & \textbf{Response} \\
\hline
\endhead

/api/auth/login & POST & \{username, password\} & \{token, user\_info\} \\
\hline
/api/auth/logout & POST & Authorization header & \{success\} \\
\hline
/api/scans/ & GET & Query params (pagination) & \{scans: [...], count\} \\
\hline
/api/scans/ & POST & FormData with MRI files & \{scan\_id, status\} \\
\hline
/api/scans/\{id\}/ & GET & scan\_id & \{scan\_details\} \\
\hline
/api/predict/ & POST & \{scan\_id\} & \{job\_id, status\} \\
\hline
/api/predict/\{id\}/ & GET & job\_id & \{status, result, progress\} \\
\hline
/api/reports/\{id\}/ & GET & report\_id & \{report\_text, json\_descriptor\} \\
\hline
/api/reports/\{id\}/review & POST & \{feedback, approved\} & \{success\} \\
\hline
/api/visualize/\{id\}/ & GET & seg\_id & \{3d\_model\_url\} \\
\hline
/api/export/\{id\}/pdf & GET & report\_id & PDF file download \\
\hline
\end{longtable}

\subsubsection{API Sequence Diagram}

\begin{figure}[H]
\centering
\includegraphics[width=\textwidth]{diagrams/api_sequence.png}
\caption{API Request-Response Sequence}
\label{fig:api_seq}
\end{figure}

\textit{Note: Generate diagram using PlantUML code provided in diagrams/api\_sequence.puml}

\subsection{User Interface Design}

\subsubsection{Dashboard Layout}

\begin{figure}[H]
\centering
\includegraphics[width=\textwidth]{diagrams/ui_dashboard.png}
\caption{Dashboard Interface Wireframe}
\label{fig:ui_dashboard}
\end{figure}

\textit{Note: Generate diagram using PlantUML code provided in diagrams/ui\_dashboard.puml}

\subsubsection{Upload Interface}

Key features:
\begin{itemize}
    \item Drag-and-drop file upload
    \item Multi-file selection for 4 modalities (T1, T1ce, T2, FLAIR)
    \item File validation (NIfTI format, size limits)
    \item Progress indicator
    \item Patient metadata form
\end{itemize}

\subsubsection{Results Visualization Interface}

Components:
\begin{itemize}
    \item 3D interactive viewer (Three.js/VTK.js)
    \item Multi-planar reconstruction (MPR) views
    \item Segmentation overlay toggle
    \item Volume measurements display
    \item Atlas mapping visualization
    \item Report text panel
\end{itemize}

\subsection{External Interface Requirements}

\subsubsection{Hardware Interfaces}
\begin{itemize}
    \item \textbf{GPU:} CUDA-compatible GPU (minimum 8GB VRAM) for inference
    \item \textbf{Storage:} Network-attached storage (NAS) or cloud storage (S3/MinIO) for MRI data
\end{itemize}

\subsubsection{Software Interfaces}
\begin{itemize}
    \item \textbf{PACS Integration:} DICOM receiver for automated MRI import
    \item \textbf{HIS/RIS:} HL7 FHIR API for patient demographics
    \item \textbf{Cloud Services:} AWS/GCP for scalable deployment
\end{itemize}

\subsubsection{Communication Interfaces}
\begin{itemize}
    \item \textbf{Protocol:} HTTPS (TLS 1.3)
    \item \textbf{Data Format:} JSON for API, NIfTI for medical images
    \item \textbf{WebSocket:} Real-time progress updates during inference
\end{itemize}

\newpage

% 6. Security Design
\section{Security Design}

\subsection{Authentication and Authorization}

\subsubsection{Authentication Flow}

\begin{figure}[H]
\centering
\includegraphics[width=\textwidth]{diagrams/auth_flow.png}
\caption{Authentication Flow Diagram}
\label{fig:auth_flow}
\end{figure}

\textit{Note: Generate diagram using PlantUML code provided in diagrams/auth\_flow.puml}

\subsubsection{JWT Token Structure}
\begin{lstlisting}[language=json]
{
  "header": {
    "alg": "RS256",
    "typ": "JWT"
  },
  "payload": {
    "user_id": "uuid",
    "username": "string",
    "role": "admin|radiologist|researcher",
    "exp": "timestamp",
    "iat": "timestamp"
  },
  "signature": "..."
}
\end{lstlisting}

\subsection{Role-Based Access Control (RBAC)}

\begin{table}[H]
\centering
\begin{tabularx}{\textwidth}{|l|X|X|X|}
\hline
\textbf{Resource} & \textbf{Admin} & \textbf{Radiologist} & \textbf{Researcher} \\
\hline
Upload Scans & ✓ & ✓ & ✗ \\
\hline
View Scans & ✓ & ✓ (own only) & ✓ (anonymized) \\
\hline
Run Inference & ✓ & ✓ & ✓ \\
\hline
Review Reports & ✓ & ✓ & ✗ \\
\hline
Export Data & ✓ & ✓ & ✓ (anonymized) \\
\hline
Manage Users & ✓ & ✗ & ✗ \\
\hline
Fine-tune Models & ✓ & ✗ & ✓ \\
\hline
\end{tabularx}
\end{table}

\subsection{Data Security}

\subsubsection{Encryption}
\begin{itemize}
    \item \textbf{At Rest:} AES-256 encryption for stored MRI files
    \item \textbf{In Transit:} TLS 1.3 for all API communications
    \item \textbf{Database:} Transparent Data Encryption (TDE) for PostgreSQL
\end{itemize}

\subsubsection{Data Anonymization}
\begin{equation}
PatientCode = HMAC-SHA256(PatientID, SecretKey)
\end{equation}

All personally identifiable information (PII) is removed before research access.

\subsection{HIPAA Compliance}

Key measures:
\begin{itemize}
    \item Audit logging for all data access
    \item Automatic session timeout (30 minutes)
    \item Password complexity requirements
    \item Multi-factor authentication (optional)
    \item Data retention policies
    \item Secure backup and disaster recovery
\end{itemize}

\newpage

% 7. Deployment Design
\section{Deployment Design}

\subsection{Deployment Architecture}

\begin{figure}[H]
\centering
\includegraphics[width=\textwidth]{diagrams/deployment_diagram.png}
\caption{Deployment Architecture Diagram}
\label{fig:deployment}
\end{figure}

\textit{Note: Generate diagram using PlantUML code provided in diagrams/deployment\_diagram.puml}

\subsection{Container Architecture}

\subsubsection{Docker Services}

\begin{table}[H]
\centering
\begin{tabularx}{\textwidth}{|l|l|X|l|}
\hline
\textbf{Service} & \textbf{Image} & \textbf{Purpose} & \textbf{Ports} \\
\hline
nginx & nginx:alpine & Reverse proxy, static files & 80, 443 \\
\hline
backend & django:custom & API server & 8000 \\
\hline
celery-worker & django:custom & Async tasks (inference) & - \\
\hline
celery-beat & django:custom & Scheduled tasks & - \\
\hline
redis & redis:7-alpine & Message broker, cache & 6379 \\
\hline
postgres & postgres:14 & Database & 5432 \\
\hline
minio & minio/minio & Object storage & 9000, 9001 \\
\hline
\end{tabularx}
\end{table}

\subsection{Scalability Design}

\subsubsection{Horizontal Scaling}
\begin{itemize}
    \item \textbf{API Server:} Load-balanced multiple instances
    \item \textbf{Celery Workers:} Auto-scaling based on queue length
    \item \textbf{Database:} Read replicas for query load distribution
\end{itemize}

\subsubsection{Vertical Scaling}
\begin{itemize}
    \item \textbf{GPU Workers:} Dedicated high-memory instances for model inference
    \item \textbf{Database:} SSD storage with sufficient IOPS
\end{itemize}

\subsection{Performance Requirements}

\begin{table}[H]
\centering
\begin{tabularx}{\textwidth}{|l|X|l|}
\hline
\textbf{Metric} & \textbf{Description} & \textbf{Target} \\
\hline
API Response Time & Time for API endpoints to respond & < 200ms \\
\hline
Segmentation Time & Time to segment single case & < 30s \\
\hline
Report Generation Time & Time to generate clinical report & < 10s \\
\hline
Concurrent Users & Number of simultaneous users & 50+ \\
\hline
Throughput & Cases processed per hour & 100+ \\
\hline
Availability & System uptime & 99.5\% \\
\hline
\end{tabularx}
\end{table}

\subsection{Backup and Recovery}

\subsubsection{Backup Strategy}
\begin{itemize}
    \item \textbf{Database:} Automated daily backups with 30-day retention
    \item \textbf{MRI Files:} Replicated to secondary storage (S3 Glacier)
    \item \textbf{Model Checkpoints:} Version-controlled with DVC (Data Version Control)
\end{itemize}

\subsubsection{Disaster Recovery}
\begin{itemize}
    \item \textbf{RTO (Recovery Time Objective):} 4 hours
    \item \textbf{RPO (Recovery Point Objective):} 24 hours
    \item \textbf{Failover:} Multi-region deployment for critical components
\end{itemize}

\newpage

% 8. Use Case Diagrams
\section{Use Case Specifications}

\subsection{Use Case Diagram}

\begin{figure}[H]
\centering
\includegraphics[width=\textwidth]{diagrams/use_case_diagram.png}
\caption{System Use Case Diagram}
\label{fig:usecase}
\end{figure}

\textit{Note: Generate diagram using PlantUML code provided in diagrams/use\_case\_diagram.puml}

\subsection{Use Case: Upload MRI Scan}

\begin{table}[H]
\centering
\begin{tabularx}{\textwidth}{|l|X|}
\hline
\textbf{ID} & UC-001 \\
\hline
\textbf{Name} & Upload MRI Scan \\
\hline
\textbf{Actor} & Radiologist \\
\hline
\textbf{Precondition} & User is authenticated \\
\hline
\textbf{Postcondition} & MRI scan is stored and ready for processing \\
\hline
\textbf{Main Flow} & 
1. User navigates to upload page \newline
2. User selects 4 MRI modality files (T1, T1ce, T2, FLAIR) \newline
3. User enters patient metadata (age, sex, scan date) \newline
4. System validates file formats and sizes \newline
5. System uploads files to object storage \newline
6. System creates scan record in database \newline
7. System displays success message \\
\hline
\textbf{Alternative Flow} & 
4a. File validation fails \newline
\quad 4a1. System displays error message \newline
\quad 4a2. User corrects the issue \newline
\quad 4a3. Return to step 2 \\
\hline
\end{tabularx}
\end{table}

\subsection{Use Case: Run Segmentation Analysis}

\begin{table}[H]
\centering
\begin{tabularx}{\textwidth}{|l|X|}
\hline
\textbf{ID} & UC-002 \\
\hline
\textbf{Name} & Run Segmentation Analysis \\
\hline
\textbf{Actor} & Radiologist, Researcher \\
\hline
\textbf{Precondition} & MRI scan is uploaded \\
\hline
\textbf{Postcondition} & Segmentation result is available \\
\hline
\textbf{Main Flow} & 
1. User selects an uploaded scan \newline
2. User clicks "Run Analysis" button \newline
3. System creates asynchronous job \newline
4. System preprocesses MRI data \newline
5. System runs MoME+ inference \newline
6. System performs atlas mapping \newline
7. System generates JSON descriptor \newline
8. System stores segmentation result \newline
9. System notifies user of completion \\
\hline
\textbf{Alternative Flow} & 
5a. Inference fails due to GPU error \newline
\quad 5a1. System retries on different worker \newline
\quad 5a2. If retry fails, notify user \\
\hline
\end{tabularx}
\end{table}

\subsection{Use Case: Generate Clinical Report}

\begin{table}[H]
\centering
\begin{tabularx}{\textwidth}{|l|X|}
\hline
\textbf{ID} & UC-003 \\
\hline
\textbf{Name} & Generate Clinical Report \\
\hline
\textbf{Actor} & Radiologist \\
\hline
\textbf{Precondition} & Segmentation analysis is complete \\
\hline
\textbf{Postcondition} & Clinical report is generated and displayed \\
\hline
\textbf{Main Flow} & 
1. User views segmentation result \newline
2. User clicks "Generate Report" button \newline
3. System retrieves JSON descriptor \newline
4. System sends descriptor to LLM model \newline
5. LLM generates clinical report text \newline
6. System performs consistency validation \newline
7. System stores report in database \newline
8. System displays report to user \\
\hline
\textbf{Alternative Flow} & 
6a. Consistency validation fails \newline
\quad 6a1. System flags report for manual review \newline
\quad 6a2. System displays warning to user \\
\hline
\end{tabularx}
\end{table}

\subsection{Use Case: Review and Approve Report}

\begin{table}[H]
\centering
\begin{tabularx}{\textwidth}{|l|X|}
\hline
\textbf{ID} & UC-004 \\
\hline
\textbf{Name} & Review and Approve Report \\
\hline
\textbf{Actor} & Radiologist \\
\hline
\textbf{Precondition} & Clinical report is generated \\
\hline
\textbf{Postcondition} & Report is marked as reviewed with feedback \\
\hline
\textbf{Main Flow} & 
1. User reads generated clinical report \newline
2. User reviews segmentation visualization \newline
3. User provides feedback (comments, corrections) \newline
4. User marks report as approved or rejected \newline
5. System updates report status \newline
6. System stores feedback for continual learning \\
\hline
\end{tabularx}
\end{table}

\newpage

% 9. Appendices
\section{Appendices}

\subsection{Appendix A: PlantUML Diagram Codes}

All PlantUML codes are provided in the \texttt{diagrams/} directory with the following files:

\begin{itemize}
    \item \texttt{system\_architecture.puml} - System architecture diagram
    \item \texttt{component\_diagram.puml} - Component relationships
    \item \texttt{erd\_diagram.puml} - Entity-relationship diagram
    \item \texttt{data\_flow\_diagram.puml} - Data flow visualization
    \item \texttt{preprocessing\_class.puml} - Preprocessing classes
    \item \texttt{mome\_architecture.puml} - MoME+ model architecture
    \item \texttt{atlas\_mapping\_sequence.puml} - Atlas mapping sequence
    \item \texttt{report\_generation\_activity.puml} - Report generation workflow
    \item \texttt{scan\_state\_diagram.puml} - Scan processing states
    \item \texttt{api\_sequence.puml} - API interaction sequence
    \item \texttt{ui\_dashboard.puml} - Dashboard wireframe
    \item \texttt{auth\_flow.puml} - Authentication flow
    \item \texttt{deployment\_diagram.puml} - Deployment architecture
    \item \texttt{use\_case\_diagram.puml} - System use cases
\end{itemize}

\subsection{Appendix B: Model Performance Metrics}

\begin{table}[H]
\centering
\begin{tabularx}{\textwidth}{|l|X|c|c|c|}
\hline
\textbf{Dataset} & \textbf{Task} & \textbf{Dice (WT)} & \textbf{Dice (TC)} & \textbf{Dice (ET)} \\
\hline
BraTS 2021 & Glioma Segmentation & 0.912 & 0.867 & 0.823 \\
\hline
OASIS & Alzheimer's Detection & 0.895 & 0.851 & - \\
\hline
ISLES & Stroke Lesion & 0.878 & - & - \\
\hline
\end{tabularx}
\caption{Expected Model Performance (Baseline)}
\end{table}

\subsection{Appendix C: Hardware Requirements}

\subsubsection{Development Environment}
\begin{itemize}
    \item \textbf{CPU:} Intel i7/AMD Ryzen 7 or better
    \item \textbf{RAM:} 32GB minimum, 64GB recommended
    \item \textbf{GPU:} NVIDIA RTX 3080 (10GB) or better
    \item \textbf{Storage:} 500GB SSD for code, 2TB HDD for datasets
\end{itemize}

\subsubsection{Production Environment}
\begin{itemize}
    \item \textbf{Inference Server:} NVIDIA A100 (40GB) or V100 (32GB)
    \item \textbf{API Server:} 8-core CPU, 16GB RAM
    \item \textbf{Database Server:} 4-core CPU, 16GB RAM, SSD storage
    \item \textbf{Storage:} 10TB+ for MRI archive, scalable object storage
\end{itemize}

\subsection{Appendix D: Third-Party Libraries}

\begin{table}[H]
\centering
\begin{tabularx}{\textwidth}{|l|l|X|}
\hline
\textbf{Library} & \textbf{Version} & \textbf{License} \\
\hline
PyTorch & 2.0+ & BSD-3-Clause \\
\hline
MONAI & 1.3+ & Apache 2.0 \\
\hline
Django & 4.2+ & BSD \\
\hline
Django REST Framework & 3.14+ & BSD \\
\hline
Celery & 5.3+ & BSD \\
\hline
nibabel & 5.1+ & MIT \\
\hline
nilearn & 0.10+ & BSD \\
\hline
transformers & 4.30+ & Apache 2.0 \\
\hline
React & 18+ & MIT \\
\hline
Three.js & 0.150+ & MIT \\
\hline
\end{tabularx}
\end{table}

\subsection{Appendix E: Glossary of Medical Terms}

\begin{table}[H]
\centering
\begin{tabularx}{\textwidth}{|l|X|}
\hline
\textbf{Term} & \textbf{Definition} \\
\hline
Glioma & A type of tumor that occurs in the brain and spinal cord \\
\hline
MRI & Magnetic Resonance Imaging - non-invasive imaging technique \\
\hline
FLAIR & Fluid-Attenuated Inversion Recovery MRI sequence useful for detecting lesions \\
\hline
Segmentation & Process of partitioning an image into meaningful regions \\
\hline
Necrosis & Death of tissue due to disease or injury \\
\hline
Edema & Swelling caused by excess fluid trapped in body tissues \\
\hline
Enhancing & Areas that show increased signal after contrast agent administration \\
\hline
Atlas & Standardized 3D brain model used for spatial normalization \\
\hline
Voxel & 3D pixel (volume element) in medical imaging \\
\hline
\end{tabularx}
\end{table}

\end{document}
